\section{The bar complex as Swiss-cheese algebra}
\label{sec:foundations}
\label{sec:operads-factorization}

\subsection{From Volume~I to three dimensions}

The bar complex $\barB^{\mathrm{ch}}(\cA)$
of a chiral algebra~$\cA$ on a curve~$X$ is the tensor coalgebra $T^c(s^{-1}\bar{\cA})$, equipped with two structures:
\begin{enumerate}[label=(\roman*)]
\item a \emph{differential} $d_{\barB}$, built from OPE
  residues along collision divisors in $\FM_k(X)$, encoding the
  chiral product;
\item a \emph{coassociative coproduct}
  $\Delta[a_1|\cdots|a_n]
  = \sum_{i=0}^{n} [a_1|\cdots|a_i]
    \otimes [a_{i+1}|\cdots|a_n]$,
  the ordered deconcatenation of tensor factors.
\end{enumerate}
The differential lives on $\FM_k(\C)$.  The coproduct lives on
$\Conf_k(\R)$: the ordered splitting of a bar element into two
sub-sequences is the algebraic expression of cutting an interval at
a point.  Together, a bar element of degree~$k$ is parametrized by
$\FM_k(\C) \times \Conf_k(\R)$ --- the product of holomorphic and
topological configuration spaces.

This product is the operadic fingerprint of a three-dimensional
holomorphic--topological (HT) field theory on $\C_z \times \R_t$:
local observables factorize holomorphically in~$z$ and
associatively in~$t$.  The governing structure is the two-colored
Swiss-cheese operad $\SCchtop$ with operation spaces
$\FM_k(\C) \times E_1(m)$.

\subsection{The correspondence structure}
\label{subsec:correspondence-structure}

The analogy between the bar complex and the Steinberg variety is not
merely heuristic.  The bar complex $\barB^{\mathrm{ch}}(\cA) =
T^c(s^{-1}\bar{\cA})$ on $\FM_k(\C) \times \Conf_k(\R)$ is a
correspondence in the precise sense that its differential arises from
convolution and its coproduct arises from the diagonal, exactly as
the Hecke algebra arises from convolution on $\widetilde{\mathcal N}
\times_{\mathfrak g} \widetilde{\mathcal N}$ and restriction along
the groupoid diagonal.

\begin{proposition}[Bar complex as correspondence algebra;
  \ClaimStatusProvedHere]
\label{prop:bar-correspondence}
Let\/ $\cA$ be a chiral algebra on\/ $\C$ with bar complex\/
$\barB^{\mathrm{ch}}(\cA) = T^c(s^{-1}\bar{\cA})$.  The bar data
$(d_{\barB}, \Delta)$ admit the following correspondence-geometric
description:
\begin{enumerate}[label=\textup{(\alph*)}]
\item \textbf{Convolution differential.}
  Consider the triple correspondence
  \[
    \FM_k(\C)
    \xleftarrow{\;p_{12}\;}
    \FM_k(\C) \times_{\partial} \FM_k(\C)
    \xrightarrow{\;p_{13}\;}
    \FM_{k-1}(\C),
  \]
  where the fiber product is taken over the collision divisors\/
  $D_S \subset \partial\FM_k(\C)$.  The bar differential\/
  $d_{\barB}$ is the convolution: the pushforward\/
  $p_{13,*}$ of the pullback\/ $p_{12}^* \boxtimes p_{23}^*$ along
  this triple correspondence.  Each summand in\/ $d_{\barB}$
  is an OPE residue along a collision stratum, which is precisely
  the external product on\/ $\FM_2 \times \FM_2$, restricted to\/
  $\FM_3$ via the boundary inclusion\/ $D_S \hookrightarrow \FM_k$,
  then pushed forward by collapse.

\item \textbf{Diagonal coproduct.}
  The deconcatenation coproduct\/
  $\Delta\colon T^c(s^{-1}\bar{\cA}) \to
  T^c(s^{-1}\bar{\cA}) \otimes T^c(s^{-1}\bar{\cA})$
  is the pullback\/ $\delta^*$ along the diagonal\/
  $\Conf_k(\R) \to \Conf_k(\R) \times \Conf_k(\R)$.  Because the
  ordered configuration space\/ $\Conf_k^{<}(\R)$ is the open
  simplex\/ $\{t_1 < \cdots < t_k\}$, hence contractible,
  the only coassociative coalgebra structure on the cofree
  conilpotent coalgebra\/ $T^c(s^{-1}\bar{\cA})$ compatible with
  the cogenerators is deconcatenation --- the diagonal of
  the groupoid of ordered interval decompositions.

\item \textbf{Coderivation = groupoid compatibility.}
  The identity
  \[
    d_{\barB} \circ \Delta
    \;=\;
    (d_{\barB} \otimes \mathrm{id}
     + \mathrm{id} \otimes d_{\barB})
    \circ \Delta
  \]
  states that the convolution differential is a coderivation of the
  diagonal coproduct: the diagonal is a groupoid morphism compatible
  with convolution.  This is the content of the bar construction:
  giving a coderivation of the cofree coalgebra is equivalent to
  giving a collection of maps\/
  $\{m_k \colon \bar{\cA}^{\otimes k} \to \cA\}_{k \geq 1}$,
  which are the\/ $A_\infty$ operations of~$\cA$.
\end{enumerate}
\end{proposition}

\begin{proof}
\textbf{Part~(a).}
The bar differential $d_{\barB}$ is the coderivation of
$T^c(s^{-1}\bar{\cA})$ determined by the collection of
$A_\infty$ operations $\{m_k\}_{k \geq 1}$.  Its $m_2$-component
acts on a bar element $[a_1|\cdots|a_n]$ by applying $m_2$ to each
\emph{consecutive} pair:
\[
  d_{\barB}^{(2)}[a_1|\cdots|a_n]
  = \sum_{i=1}^{n-1}
    \pm\, [a_1|\cdots|m_2(a_i,a_{i+1})|\cdots|a_n],
\]
where $m_2(a_i,a_{i+1}) = \Res_{z_i = z_{i+1}} \mu(a_i,a_{i+1})$
is the OPE residue along the collision divisor
$D_{\{i,i+1\}} \subset \partial\FM_n(\C)$.  This is the
convolution recipe: form the external product of $a_i$ and
$a_{i+1}$ on the two-point space $\FM_2(\C)$, restrict to the
fiber product over the boundary stratum
$D_{\{i,i+1\}} \hookrightarrow \FM_n(\C)$, and push forward along
the collapse map $\FM_n(\C) \to \FM_{n-1}(\C)$ that merges
$z_i$ and $z_{i+1}$.  The higher components $d_{\barB}^{(k)}$
for $k \geq 3$, built from the operations $m_k$, arise from
deeper boundary strata --- simultaneous collisions of $k$
consecutive points --- which correspond to iterated convolutions
along higher fiber products.

\textbf{Part~(b).}
The ordered configuration space $\Conf_k^{<}(\R) = \{t_1 < \cdots
< t_k\} \subset \R^k$ is an open simplex, hence contractible.
On a contractible space, a locally constant cosheaf is determined by
its value at a point.  The cofree conilpotent coalgebra
$T^c(V)$ on a graded vector space $V$ admits a unique coassociative
coproduct compatible with the coaugmentation filtration, namely
deconcatenation (this is the defining universal property of the
cofree coalgebra; see e.g.\ Loday--Vallette~\cite{LV12},
Proposition~1.2.9).  Geometrically, the deconcatenation
$[a_1|\cdots|a_n] \mapsto \sum_{i=0}^{n} [a_1|\cdots|a_i] \otimes
[a_{i+1}|\cdots|a_n]$ is the diagonal of the self-intersection
groupoid: cutting the interval $[t_1,t_n]$ at $t_i$ gives two
sub-intervals whose ordered configurations assemble into the two
tensor factors.

\textbf{Part~(c).}
A linear map $d\colon T^c(V) \to T^c(V)$ of degree~$+1$ is a
coderivation of the deconcatenation coproduct if and only if it is
determined by its components $\pi \circ d|_{V^{\otimes k}} \colon
V^{\otimes k} \to V$, where $\pi\colon T^c(V) \to V$ is the
cogenerator projection (Loday--Vallette~\cite{LV12},
Proposition~1.3.3).  These components are precisely the
$A_\infty$ operations $\{m_k\}_{k \geq 1}$.  The condition
$d_{\barB}^{\,2} = 0$ then encodes the $A_\infty$ relations
$\sum_{p+q+r=n} m_{p+1+r} \circ (\mathrm{id}^{\otimes p} \otimes
m_q \otimes \mathrm{id}^{\otimes r}) = 0$, which are exactly the
Stasheff identities.  This is Vol~I, Theorem~A: the bar differential
is a coderivation, and the bar-cobar adjunction is the translation
between coderivations of the cofree coalgebra and $A_\infty$
structure maps.
\end{proof}

\begin{corollary}[$\R$-factorization from the Lagrangian embedding;
  \ClaimStatusProvedHere]
\label{cor:R-factorization-lagrangian}
The topological coproduct\/ $\Delta$ carries no independent data
beyond the holomorphic differential.  Precisely: the
$\R$-factorization is determined by the Lagrangian embedding\/
$\mathcal L \hookrightarrow \mathcal M$ as the unique
coassociative coproduct on the cofree coalgebra\/
$T^c(s^{-1}\bar{\cA})$.  Geometrically, $\Delta$ is the diagonal of
the self-intersection groupoid\/
$\mathcal L \times_{\mathcal M} \mathcal L$, the desuspension\/
$s^{-1}$ is the cotangent shift\/ $T[-1]$ forced by the Lagrangian
condition, and the cofree coalgebra is the symmetric algebra on
conormal fibers of the embedding.
\end{corollary}

\begin{proof}
By Proposition~\ref{prop:bar-correspondence}(b), the
deconcatenation coproduct is the unique coassociative coproduct on
$T^c(s^{-1}\bar{\cA})$ compatible with the cogenerators.  It is
therefore determined by the cogenerating space
$s^{-1}\bar{\cA}$, which is the data of the holomorphic chiral
algebra~$\cA$ (with degree shift).  The Lagrangian interpretation
follows: the cofree coalgebra $T^c(N^*[-1])$ on the shifted
conormal bundle of an embedding $\mathcal L \hookrightarrow
\mathcal M$ carries a canonical deconcatenation coproduct, and the
self-intersection $\mathcal L \times_{\mathcal M} \mathcal L$ is
presented by this coalgebra with diagonal given by deconcatenation.
Here $\mathcal L$ is the support of the chiral algebra
(the Ran space of~$\C$), $\mathcal M$ is the ambient moduli, and
the Lagrangian condition forces the $(-1)$-shift that produces the
bar desuspension $s^{-1}$.
\end{proof}

The central claim of this volume: \emph{the bar complex of
Volume~I, equipped with its coproduct, is a coalgebra
over~$\SCchtop$.}  The bar differential is holomorphic
factorization; the bar coproduct is topological factorization; the
no-open-to-closed rule is the physical directionality that bulk
interactions restrict to boundaries but not conversely.  By operadic Koszul self-duality
(\S\ref{subsec:dg-yangian-operadic}), $\cA^!$ inherits an
$\SCchtop$-algebra structure---a dg-shifted Yangian
(Theorem~\ref{thm:Koszul_dual_Yangian}).

A deeper fact (Volume~I,
Theorem~\ref*{thm:bar-swiss-cheese} of
the $E_n$-Koszul duality chapter) sharpens
this claim: the $\R$-factorization is \emph{determined}
by the modular homotopy type of~$\cA$.  The
bar construction $T^c(s^{-1}\bar{\cA})$ is the cofree conilpotent
coalgebra on $s^{-1}\bar{\cA}$; its deconcatenation coproduct is
the unique coassociative coproduct compatible with the
cogenerators; and the bar-cobar adjunction (Volume~I, Theorem~A)
together with bar-cobar inversion (Volume~I, Theorem~B) guarantee
that this tensor coalgebra presentation is canonical up to
contractible ambiguity.  At genus~$g \geq 1$, the scalarity of the
modular characteristic~$\kappa(\cA)$ (Volume~I, Theorem~D) ensures
that the curvature $\dfib^2 = \kappa(\cA) \cdot \omega_g$ is
central in the coalgebra, so the coderivation property persists and
the product structure $\FM_k(\C) \times \Conf_k(\R)$ survives.

\begin{principle}[$\R$-factorization as shadow]
\label{princ:R-shadow}
The $\R$-factorization is the shadow of the modular homotopy type
projected onto the topological direction.  Given holomorphic
factorization on $\operatorname{Ran}(X)$, the bar-cobar
adjunction produces a canonical tensor coalgebra model; the
$\Eone$-coalgebra structure is automatic.  The bar construction
is the \emph{universal boundary condition}: the cofree resolution
of~$\cA$ viewed from the~$\R$-direction.  The
Swiss-cheese structure of this volume is not an ansatz but a theorem.
\end{principle}


\subsection{The two–colored Swiss–cheese operad}
\label{subsec:SC-operad}

\begin{definition}[Colors]
\label{def:SC-colors}
Consider the set of two colors $\mathbf{Col}=\{\mathsf{ch},\mathsf{top}\}$, thought of as the
\emph{closed} (holomorphic) and \emph{open} (topological) colors respectively.
\end{definition}

\begin{definition}[Spaces of operations]
\label{def:SC-operations}
Define the topological 2–colored operad
$\mathsf{SC}^{\mathrm{ch,top}}$ as follows.
\begin{enumerate}[label=(\roman*)]
\item \textbf{Closed outputs (holomorphic):} For closed output and only closed inputs,
  set
  \[
     \mathsf{SC}^{\mathrm{ch,top}}\!\bigl( (\underbrace{\mathsf{ch},\ldots,\mathsf{ch}}_{k});\,\mathsf{ch} \bigr)
     \;:=\; \FM_k(\mathbb C),
  \]
  the Fulton–MacPherson compactification of the configuration space of $k$ points in $\mathbb C$.
  If any input is open, the space of operations with closed output is \emph{empty}:
  $\mathsf{SC}^{\mathrm{ch,top}}\bigl( (\ldots,\mathsf{top},\ldots);\,\mathsf{ch} \bigr)=\varnothing$.
  By convention, $\FM_0(\C)$ and $\FM_1(\C)$ are single points, and $E_1(0) = \emptyset$ (the empty configuration).
\item \textbf{Open outputs (topological):} For open output and a mixture of open and closed inputs, set
  \[
     \mathsf{SC}^{\mathrm{ch,top}}\!\bigl( (\underbrace{\mathsf{ch},\ldots,\mathsf{ch}}_{k}, \underbrace{\mathsf{top},\ldots,\mathsf{top}}_{m});\,\mathsf{top} \bigr)
     \;:=\; \FM_k(\mathbb C)\times E_1(m) .
  \]
  This encodes that an open operation may depend on closed inputs (holomorphic insertions)
  and on open inputs (time–ordered insertions), and its structure is the product of
  holomorphic and topological composition spaces.
\end{enumerate}
\end{definition}

\begin{remark}[Why the product structure is forced]
\label{rem:product-forced-by-locality}
The product decomposition
$\FM_k(\C) \times E_1(m)$ in
Definition~\ref{def:SC-operations}(ii) is forced by locality: a
local observable that factorizes holomorphically in~$z$ and is
locally constant in~$t$ can only depend on collision data in~$\C$
(parametrized by $\FM_k(\C)$) and time-ordering in~$\R$
(parametrized by $\Conf_k^{<}(\R) \subset E_1(m)$).  No other
topological data is available.  In particular, there is no mixed
configuration space interpolating between the holomorphic and
topological directions: the kernel of each $m_k$ integral
factorizes as a meromorphic function of $z_{ij} = z_i - z_j$ times
a step function of $t_{ij} = t_i - t_j$, and this factorization is
a consequence of the equations of motion in the HT gauge, as packaged
for physical realizations by Theorem~\ref{thm:physics-bridge}.
\end{remark}

\begin{remark}[Swiss–cheese directionality]
The absence of mixed operations $(\ldots,\mathsf{top},\ldots)\to\mathsf{ch}$ ensures a
\emph{closed–output isolation}: holomorphic (“closed”) outputs cannot depend on open inputs.
This is the hallmark of Swiss–cheese operads and reflects the physical directionality:
bulk (closed) interactions can restrict to boundaries (open), but boundary interactions
cannot produce new bulk structures.
\end{remark}

\begin{definition}[Composition]
\label{def:SC-composition}
Composition in $\mathsf{SC}^{\mathrm{ch,top}}$ is induced by the strict operadic insertion
maps of $\FM_\bullet(\mathbb C)$ and $E_1$, composed \emph{componentwise} as follows.
\begin{itemize}
  \item For closed–into–closed insertions,
    \[
      \FM_k(\mathbb C)\times\prod_{i=1}^k\FM_{\ell_i}(\mathbb C)
      \;\longrightarrow\;
      \FM_{\sum_i\ell_i}(\mathbb C),
    \]
    we use the standard operadic insertion of configurations of points.
  \item For open–into–open with $k$ closed and $m$ open inputs,
    \[
      \bigl(\FM_k(\mathbb C)\times E_1(m)\bigr)\times
      \prod_{i=1}^k\FM_{\ell_i}(\mathbb C)\times\prod_{j=1}^m E_1(n_j)
      \;\longrightarrow\;
      \FM_{\sum_i\ell_i}(\mathbb C)\times E_1(\sum_j n_j) ,
    \]
    taking the product of insertion maps on $\FM_\bullet(\mathbb C)$ and on $E_1$.
  \item For closed–into–open, we allow insertion of closed operations in the closed $\FM$–slots
    inside open operations, but never the reverse (no open–into–closed).
\end{itemize}
The symmetric groups act by permuting labeled inputs of each color, and units are given by
the arity–1 elements in $\FM_1(\mathbb C)$ and~$E_1(1)$.
\end{definition}

\begin{proposition}[Well-definedness; \ClaimStatusProvedHere]
\label{prop:SC-well-defined}
The data of Definitions~\ref{def:SC-colors}–\ref{def:SC-composition}
endows $\mathsf{SC}^{\mathrm{ch,top}}$ with the structure of a
strict two–colored topological operad.
\end{proposition}

\begin{proof}
Associativity, unitality, and symmetric–group compatibility follow from the corresponding
properties of the operads $\FM_\bullet(\mathbb C)$ and~$E_1$, together with the product structure of the
mixed composition.  The "no open-to-closed" restriction is
stable under composition by construction.
\end{proof}

\begin{remark}[Homotopy-Koszulity and two-color Koszul duality]
\label{rem:two-color-koszul-duality-operadic}
\index{two-color Koszul duality!operadic input}
The homotopy-Koszulity of $\SCchtop$
(Theorem~\ref{thm:homotopy-Koszul}) is the operadic input to
the master two-color Koszul duality theorem
(Theorem~\ref{thm:two-color-master}).  Together with Koszul
self-duality (\S\ref{subsec:dg-yangian-operadic}), it ensures
that $\cA^!$ is an $\SCchtop$-algebra---a dg-shifted Yangian
(Theorem~\ref{thm:Koszul_dual_Yangian}).

The associated graded splitting
$\mathrm{gr}\;\SCchtop \cong \mathsf{P}_{\mathrm{ch}}
\amalg E_1$
(Proposition~\ref{prop:gr-chiral}) shows that the two colors
decouple at the quadratic level; the full $\SCchtop$-algebra
is a \emph{filtered deformation} of this decoupled structure.
The holomorphic weight filtration is the organizing filtration
that separates the ``easy'' (decoupled) data from the
``interacting'' data: the differentials in the associated
spectral sequence encode successively higher-order corrections,
starting from the classical $r$-matrix at $E_1$ and converging
to the full quantum $R(z)$.
\end{remark}

\begin{construction}[$E_n$ formality hierarchy as MC truncation]
\label{constr:vol2-en-formality}
The $E_n$ formality hierarchy corresponds to truncation of the MC obstruction tower, as in the $E_n$-formality/MC-truncation proposition of Volume~I:
\begin{enumerate}[label=\textup{(\roman*)}]
\item $E_\infty$-formality: all higher brackets $\ell_n^{(g)} = 0$ for $n \geq 3$ in $\gAmod$.
\item $E_2$-formality: genus-$0$ convolution algebra $\mathfrak{g}^{(0)}_\cA$ is formal (chiral Koszulness).
\item $E_1$-formality: ordered convolution algebra $\mathfrak{g}^{E_1}_\cA$ carries a formal $A_\infty$-model.
\end{enumerate}
The Swiss-cheese operad $\mathrm{SC}^{\mathrm{ch,top}}$ is homotopy-Koszul (Theorem~\ref{thm:homotopy-Koszul}), placing the present volume at the $E_1$ level.
\end{construction}

\subsection{Algebras and the Boardman–Vogt resolution}
\label{subsec:WSC-algebras}

\begin{definition}[Algebras over $\mathsf{SC}^{\mathrm{ch,top}}$]
\label{def:SC-algebra}
Let $(\mathcal C,\otimes,\mathbf 1)$ be a symmetric monoidal presentable stable dg–category
(e.g.\ chain complexes over a field $k$ of characteristic zero).  An algebra over the operad
$C_\ast(\mathsf{SC}^{\mathrm{ch,top}})$ in $\mathcal C$ is a pair of objects
$(A_{\mathsf{ch}},A_{\mathsf{top}})\in\mathcal C^2$
together with structure maps
\[
  C_\ast(\FM_k(\mathbb C))\otimes A_{\mathsf{ch}}^{\otimes k}\longrightarrow A_{\mathsf{ch}} ,
  \qquad
  C_\ast(\FM_k(\mathbb C)\times E_1(m))\otimes
  A_{\mathsf{ch}}^{\otimes k}\otimes A_{\mathsf{top}}^{\otimes m}\longrightarrow A_{\mathsf{top}} ,
\]
satisfying the operadic identities.
\end{definition}

\begin{definition}[Boardman–Vogt resolution]
\label{def:WSC}
Let $W(\mathsf{SC}^{\mathrm{ch,top}})$ denote the Boardman–Vogt resolution of the
topological operad $\mathsf{SC}^{\mathrm{ch,top}}$.
Its object $W(\mathsf{SC}^{\mathrm{ch,top}})$ is a cofibrant replacement in the model category of two-colored topological operads \cite{BM03,CM14}, encoding
the data of \emph{homotopy–coherent} operations and relations.
An algebra over $C_\ast\!\bigl(W(\mathsf{SC}^{\mathrm{ch,top}})\bigr)$ is called a
\emph{homotopy holomorphic–topological algebra}.
\end{definition}

\begin{remark}[Homotopy–coherent composition]
\label{rmk:W-homotopy}
The $W$--resolution introduces \emph{edge lengths} on the trees that index compositions,
tracking higher homotopies between different composition orders.
These encode $A_\infty$--type structures in the topological color
and chiral associativity up to homotopy in the closed color.
\end{remark}

\subsection{HT prefactorization algebras and Swiss–cheese algebras}
\label{subsec:prefact-vs-operad}

\begin{definition}[HT prefactorization algebra]
\label{def:HT-prefactorization}
A \emph{holomorphic–topological prefactorization algebra} on $\mathbb C\times\mathbb R$ with values
in $\mathcal C$ assigns to every sufficiently nice open set $U\subset\mathbb C\times\mathbb R$ an
object $\mathsf{Obs}(U)\in\mathcal C$ together with
\begin{enumerate}[label=(\alph*)]
\item \emph{multiplicative structure} for disjoint unions
  $U=\bigsqcup_i U_i$:
  \(
    \bigotimes_i \mathsf{Obs}(U_i)\to \mathsf{Obs}(U),
  \)
\item \emph{restriction structure} for inclusions $U\subset V$:
  \(
    \mathsf{Obs}(U)\to \mathsf{Obs}(V),
  \)
\end{enumerate}
which is (i) holomorphic in the $\mathbb C$–direction, with meromorphic singularities (OPE poles) arising from the collision of operator insertions,
(ii) locally constant in the $\mathbb R$–direction (topological factorization),
and (iii) compatible with iterated decomposition by rectangles $D\times I$.
\end{definition}

\begin{theorem}[Swiss--cheese recognition; \ClaimStatusProvedHere]
\label{thm:recognition-foundations}
\label{thm:recognition}
Let $\mathcal C$ be a symmetric monoidal presentable stable
dg--category over a field $k$ of characteristic zero.
The following data are equivalent:
\begin{enumerate}[label=(\roman*)]
\item A homotopy holomorphic--topological algebra over
  $C_\ast\!\bigl(W(\mathsf{SC}^{\mathrm{ch,top}})\bigr)$ in
  $\mathcal C$;
\item A holomorphic--topological prefactorization algebra
  $\mathsf{Obs}$ on~$\mathbb C\times\mathbb R$ with values in
  $\mathcal C$, satisfying:
  \begin{enumerate}[label=(\alph*)]
    \item \emph{Holomorphic factorization along $\mathbb C$:} the
      assignment $D\mapsto\mathsf{Obs}(D\times I)$ for discs
      $D\subset\mathbb C$ forms a prefactorization algebra
      on~$\mathbb C$ governed by
      $C_\ast(\FM_\bullet(\mathbb C))$, with meromorphic OPE
      singularities along collision diagonals;
    \item \emph{Local constancy along $\mathbb R$:} the assignment
      $I\mapsto\mathsf{Obs}(D\times I)$ for intervals
      $I\subset\mathbb R$ is locally constant (restriction along
      inclusions of intervals is a quasi-isomorphism), yielding an
      $E_1$--algebra structure;
    \item \emph{Color compatibility:} the closed--color
      (holomorphic) factorization acts on the open--color
      (topological) factorization via the product structure
      $\FM_k(\mathbb C)\times E_1(m)$, but no open--to--closed
      operations exist ($\mathsf{SC}^{\mathrm{ch,top}}(\ldots,
      \mathsf{top},\ldots;\mathsf{ch})=\varnothing$).
  \end{enumerate}
\end{enumerate}
\end{theorem}

\begin{proof}
The proof assembles three independently established recognition
results via the product decomposition of the operation spaces of
$\mathsf{SC}^{\mathrm{ch,top}}$.  The key technical input is the
Weiss cosheaf descent lemma
(Lemma~\ref{lem:product-weiss-descent}), which shows that HT
prefactorization algebras on rectangles satisfy descent and hence
determine operadic algebra structures.

\medskip\noindent
\textbf{Step~1: Closed-color recognition (holomorphic).}
Costello--Gwilliam~\cite{CG17} prove that the $(\infty,1)$-category
of holomorphic prefactorization algebras on~$\mathbb C$ with
meromorphic OPE singularities is equivalent to the category of
algebras over the chains
$C_\ast(\FM_\bullet(\mathbb C); \mathscr{L}_{\log})$, where
$\mathscr{L}_{\log}$ is the logarithmic local system encoding OPE
poles.  Since the closed-color component of
$\mathsf{SC}^{\mathrm{ch,top}}$ is exactly $\FM_k(\mathbb C)$
(Definition~\ref{def:SC-operations}(i)), condition~(ii)(a) is
equivalent to giving the closed-color component of a
$C_\ast(W(\mathsf{SC}^{\mathrm{ch,top}}))$-algebra structure on
$A_{\mathsf{ch}}$.

\medskip\noindent
\textbf{Step~2: Open-color recognition (topological).}
The Ayala--Francis recognition theorem~\cite{AF15} identifies
locally constant factorization algebras on~$\mathbb R$ satisfying
Weiss cosheaf descent with
$E_1$-algebras (associative algebras up to coherent homotopy).
Condition~(ii)(b) is therefore equivalent to giving the open-color
component of the algebra structure on $A_{\mathsf{top}}$.
The descent condition is verified by
Lemma~\ref{lem:product-weiss-descent}, Step~3: locally constant
prefactorization algebras on~$\R$ automatically satisfy Weiss descent
because the \v{C}ech nerve of any Weiss cover by intervals is
contractible and all restriction maps are quasi-isomorphisms.

\medskip\noindent
\textbf{Step~3: Two-colored assembly via the product decomposition.}
The operation spaces of
$\mathsf{SC}^{\mathrm{ch,top}}$ are \emph{products}:
\[
\mathsf{SC}^{\mathrm{ch,top}}
\bigl((\mathsf{ch}^k,\mathsf{top}^m);\mathsf{top}\bigr)
= \FM_k(\mathbb C)\times E_1(m).
\]
Moreover, the no-open-to-closed rule
($\mathsf{SC}^{\mathrm{ch,top}}(\ldots,\mathsf{top},\ldots;
\mathsf{ch})=\varnothing$) implies that the closed color is an
\emph{operadic retract}: restricting
$\mathsf{SC}^{\mathrm{ch,top}}$ to the closed color recovers
$\FM_\bullet(\mathbb C)$ as a sub-operad, and this restriction
admits a section (the inclusion of closed-only operations).  The
open-color operations see the closed color only through the product
factor $\FM_k(\mathbb C)$, never the reverse.

In the model category of $\Sigma$-cofibrant two-colored topological
operads~\cite{BM03,CM14}, the Boardman--Vogt resolution
$W(\mathsf{SC}^{\mathrm{ch,top}})$ is cofibrant.  The product
decomposition of the operation spaces is preserved by the
$W$-construction because the cubical edge-length parameters act
independently on the $\FM_k(\mathbb C)$ and $E_1(m)$ factors: a
$W$-tree with internal edges decorated by $s_e \in [0,1]$ decomposes
into its holomorphic ($\FM$-decorated) and topological
($E_1$-decorated) parts, with the same edge parameters governing
both.

Therefore, giving a
$C_\ast(W(\mathsf{SC}^{\mathrm{ch,top}}))$-algebra
$(A_{\mathsf{ch}}, A_{\mathsf{top}})$ is equivalent to giving:
\begin{itemize}
\item a $C_\ast(W(\FM_\bullet(\mathbb C)))$-algebra structure on
  $A_{\mathsf{ch}}$ (the closed-color data, governing holomorphic
  factorization);
\item a $C_\ast(W(E_1))$-algebra structure on
  $A_{\mathsf{top}}$ (the open-color data, governing topological
  factorization);
\item a compatible action of $A_{\mathsf{ch}}$ on
  $A_{\mathsf{top}}$ via the mixed operations
  $C_\ast(\FM_k(\mathbb C) \times E_1(m))$, coherent with the
  Boardman--Vogt homotopies.
\end{itemize}
Steps~1 and~2 identify the first two items with
conditions~(ii)(a) and~(ii)(b).  Condition~(ii)(c) is exactly the
third item: the closed-color holomorphic structure acts on the
open-color topological structure via the product operation spaces,
with no open-to-closed flow.  The no-open-to-closed rule ensures
that the closed-color recognition (Step~1) is not modified by the
presence of the open color, so the two single-colored recognitions
compose without obstruction.

\medskip\noindent
\textbf{Step~4: Weiss descent for the mixed operations.}
The remaining content is the converse direction: given an HT
prefactorization algebra $\Obs$ satisfying (ii)(a)--(c), we must
produce the mixed operadic structure maps
$C_\ast(\FM_k(\C) \times E_1(m)) \otimes A_{\mathsf{ch}}^{\otimes k}
\otimes A_{\mathsf{top}}^{\otimes m} \to A_{\mathsf{top}}$.
By Lemma~\ref{lem:product-weiss-descent}, $\Obs$ satisfies Weiss
cosheaf descent on the basis of rectangles
$\{D \times I\}$.  The product Weiss descent argument
(Lemma~\ref{lem:product-weiss-descent}, Step~4) shows that the
\v{C}ech nerve of any Weiss cover of $D \times I$ by sub-rectangles
computes $\Obs(D \times I)$ via an iterated hocolimit that factors
through the holomorphic and topological directions independently.
The mixed structure maps are then read off from the factorization
structure on rectangles $D_1 \sqcup \cdots \sqcup D_k \times
I_1 \sqcup \cdots \sqcup I_m \hookrightarrow D \times I$:
the holomorphic dependence on disk positions is governed by
$C_\ast(\FM_k(\C))$ (by Step~1 and condition~(ii)(a)),
the topological dependence on interval positions is governed by
$C_\ast(E_1(m))$ (by Step~2 and condition~(ii)(b)),
and the K\"unneth quasi-isomorphism
$C_\ast(\FM_k(\C)) \otimes C_\ast(E_1(m)) \simeq
C_\ast(\FM_k(\C) \times E_1(m))$
(valid in characteristic zero with finite-type homology)
assembles these into the required mixed operation maps.
Higher coherences are supplied by the $W$-resolution as in Step~3.

\medskip\noindent
\textbf{Conclusion.}
Assembling the four steps yields the equivalence between~(i)
and~(ii): the product decomposition of the operation spaces
(Step~3), the factor-by-factor recognition (Steps~1--2), the Weiss
cosheaf descent (Step~4 via
Lemma~\ref{lem:product-weiss-descent}), and the K\"unneth assembly
of mixed operations complete the proof.\qedhere
\end{proof}

\begin{remark}[Structure of the proof]
\label{rmk:recognition-proof-structure}
The proof rests on three independently established results: the
Costello--Gwilliam holomorphic recognition~\cite{CG17}, the
Ayala--Francis topological recognition~\cite{AF15}, and the
Berger--Moerdijk / Caviglia model structure for colored
operads~\cite{BM03,CM14}.  Two structural features reduce
the two-colored recognition to the product of two single-colored
recognitions plus a compatibility datum: the
\emph{product decomposition} of the operation spaces of
$\mathsf{SC}^{\mathrm{ch,top}}$ and the \emph{retract
property} of the closed color.  The compatibility datum
is precisely condition~(ii)(c).  The Weiss cosheaf descent lemma
(Lemma~\ref{lem:product-weiss-descent}) closes the remaining gap by
showing that the prefactorization data on small rectangles determine
the full operadic structure.
\end{remark}

\begin{remark}[Recognition as correspondence geometry]
\label{rem:recognition-correspondence}
Theorem~\ref{thm:recognition} is the statement that the Lagrangian
formal self-intersection model determines the Swiss-cheese algebra: one need not
construct the operad action on $(A_{\mathsf{ch}}, A_{\mathsf{top}})$
by hand.  The holomorphic factorization (the structure sheaf of
$\FM_k(\C)$) together with the topological local constancy (the
groupoid structure of $\Conf_k(\R)$) and the product decomposition
of the operation spaces suffice --- the full $\SCchtop$-algebra
structure is inherited from the correspondence geometry of the bar
complex, exactly as the Hecke algebra action on a module is
inherited from convolution on the Steinberg variety
(Theorem~\ref{thm:steinberg-presentation}).
\end{remark}

\begin{remark}[Comparison with the BD chiral operad]
\label{rmk:BD-comparison}
There is a well–known comparison between the BD chiral operad
$\mathsf P_{\mathrm{ch}}$ (defined algebraically by D–modules on the Ran space of~$\mathbb C$)
and the chains on the configuration–space operad
$C_\ast(\FM_\bullet(\mathbb C);\mathscr L_{\log})$ with logarithmic local systems, see
\cite{HA,CG17}.
Our Swiss–cheese construction is thereby compatible with the "chiral" closed color and the
topological $E_1$ open color in the factorization–algebra setting.
\end{remark}

\subsection{Compatibility with the operadic model}
\label{subsec:operadic-compatibility}

\begin{proposition}[Compatibility with $W(\mathsf{SC}^{\mathrm{ch,top}})$;
\ClaimStatusProvedHere]
\label{prop:compatibility}
Every logarithmic $\SCchtop$-algebra (Definition~\ref{def:log-SC-algebra}) determines a
holomorphic--topological prefactorization algebra on $\C\times\R$, locally
constant on rectangles, and hence a $C_\ast\!\bigl(W(\mathsf{SC}^{\mathrm{ch,top}})\bigr)$--algebra
by Theorem~\ref{thm:recognition}.
\end{proposition}

\begin{proof}
The product decomposition $\omega_k = \omega_k^{\mathrm{hol}} \otimes \omega_k^{\mathrm{top}}$ ensures compatibility of meromorphic OPEs with topological gluings.  $W$-coherence follows from the homological algebra of integrated weight forms over $\FM_k(\C)\times E_1(m)$; the FM calculus (Theorem~\ref{thm:FM-calculus}) guarantees independence of auxiliary choices.
\end{proof}

\subsection{Configuration spaces and higher operations}
\label{subsec:FM-higher-ops}
The higher multiplications arise by integrating out interaction-point
configurations.  In the holomorphic--topological setting,
\(
m_k: A^{\otimes k}\to A((\lambda_1))\cdots((\lambda_{k-1}))
\)
arise from the pushforward of the product of propagators over $\FM_k(\mathbb C)$ times
the ordered intervals in~$E_1$.
Boundary faces correspond to degenerations where subsets of point collide or where intervals
shrink.  The Arnold–Orlik–Solomon identities for logarithmic forms on $\FM_k(\mathbb C)$
ensure that Stokes' theorem reproduces the $A_\infty$~identities satisfied by $m_k$
(see Section~\ref{sec:FM_calculus}).

\subsection{The five theorems as Lagrangian geometry}
\label{subsec:five-theorems-lagrangian}

\providecommand{\cL}{\mathcal L}
\providecommand{\cM}{\mathcal M}
\providecommand{\HH}{\operatorname{HH}}
\providecommand{\Steinb}{\mathfrak{S}}

Volume~I establishes five theorems about the categorical logarithm.
Theorem~\ref{thm:bar-is-self-intersection} --- the bar complex is the
formal bar/Koszul model for the Lagrangian self-intersection ---
reveals their geometric content.
Each theorem is a shadow of the same Lagrangian embedding
$\cL \hookrightarrow \cM$, read through different invariants.

\begin{proposition}[The five theorems and their Lagrangian-geometric shadow;
\ClaimStatusHeuristic]
\label{prop:five-theorems-geometry}
Let\/ $\cM$ be a $(-2)$-shifted symplectic derived stack and\/
$\cL \hookrightarrow \cM$ a Lagrangian, with\/
$\Steinb = \cL \times_{\cM}^h \cL$ the derived self-intersection.
Write\/ $B = \cO(\cL)$.  Then Vol~I's five theorems admit the
following geometric readings:
\begin{enumerate}[label=\textup{(\roman*)},leftmargin=2em]
\item \textbf{Theorem~A \textup{(}Bar-cobar adjunction\textup{)}.}
  The self-intersection\/ $\Steinb = \cL \times_{\cM}^h \cL$ is a
  groupoid object over~$\cL$: source and target are the two
  projections $\Steinb \rightrightarrows \cL$, composition is the
  fiber product over~$\cM$, and the unit is the diagonal
  $\cL \to \Steinb$.  Every groupoid object in a stable category
  determines an adjunction between its category of comodules\/
  \textup{(}the bar side: cofree coalgebra on the conormal
  fibers\textup{)} and its category of modules\/
  \textup{(}the cobar side: free algebra on the normal
  fibers\textup{)}.  The bar-cobar adjunction of Volume~I is
  read on this surface through the corresponding formal
  groupoid comodule-module adjunction.

\item \textbf{Theorem~B \textup{(}Koszul inversion\textup{)}.}
  On the clean locus --- where the derived intersection has no
  higher Tor --- one expects the self-intersection groupoid\/
  $\Steinb$ to collapse to its base and the bar-cobar round-trip
  to become an equivalence.  Off the clean locus, the higher-Tor
  directions persist, and the bar-cobar map should be read as a
  completion rather than an inversion theorem.

\item \textbf{Theorem~C \textup{(}Complementarity\textup{)}.}
  Two complementary Lagrangians\/ $\cL, \cL'$ in~$\cM$ intersect
  in a derived scheme\/ $\cL \times_{\cM}^h \cL'$ carrying a
  $(-1)$-shifted symplectic structure \textup{(}PTVV
  \cite{PTVV13}\textup{)}.  The complementarity potential\/
  $S_{\cA}$ of Volume~I is read geometrically as the action
  functional attached to this $(-1)$-shifted symplectic
  intersection.
  The complementary pair decomposes the tangent space\/
  $T_p\cM \simeq T_p\cL \oplus T_p\cL'$ at every intersection
  point~$p$, and this decomposition is the complementary
  Lagrangian splitting of the obstruction complex that governs the
  genus-$g$ deformation theory.

\item \textbf{Theorem~D \textup{(}Leading coefficient\textup{)}.}
  The curvature\/ $\kappa(\cA) \cdot \omega_g$ at genus\/
  $g \geq 1$ is read as the first-order deformation of the
  Lagrangian embedding\/ $\cL \hookrightarrow \cM_{\mathrm{vac}}$
  as the moduli of curves varies over\/ $\overline{\mathcal M}_g$.
  In this dictionary, the equation
  $\dfib^2 = \kappa(\cA) \cdot \omega_g$ records the failure of the
  displaced Lagrangian to remain isotropic to first order.

\item \textbf{Theorem~H \textup{(}Hochschild ring\textup{)}.}
  The HKR theorem for Lagrangian embeddings gives
  \[
    \HH^\bullet(B)
    \;\simeq\;
    \cO\bigl(T^*[-1]\cL\bigr).
  \]
  The Hochschild cohomology of the boundary algebra~$B$ is the
  ring of functions on the shifted cotangent bundle of~$\cL$,
  which supplies the shifted-cotangent formal local bulk model.
  Recovering the actual formal neighborhood of~$\cL$ in~$\cM$
  requires the additional $(-1)$-shifted closed $1$-form from the
  formal Darboux theorem.  This is the geometric content of
  Theorem~H: the derived center of the boundary fixes the
  shifted-cotangent side of the bulk reconstruction, and the BV
  bracket on bulk observables is the $(-1)$-shifted Poisson bracket
  on $T^*[-1]\cL$ induced by the ambient $(-2)$-shifted symplectic
  structure.
\end{enumerate}
\end{proposition}

\begin{proof}
This proposition is a geometric dictionary, not a new package of
global derived-stack theorems.  Item~\textup{(i)} uses the standard
groupoid structure on\/ $\cL \times_{\cM}^h \cL$ together with
Theorem~\ref{thm:bar-is-self-intersection}(ii), which identifies the
derived tensor product with its formal bar/Koszul model.  Item~\textup{(iii)}
is the PTVV geometry of intersections of Lagrangians, read through the
complementarity package of Volume~I.  Item~\textup{(v)} is the HKR
identification of Theorem~\ref{thm:bar-is-self-intersection}(iii).

Items~\textup{(ii)} and~\textup{(iv)} are intentionally weaker: they
record the expected global geometry suggested by the formal model.
Proving them as statements about actual \'etale groupoids or
Kodaira--Spencer classes of global derived stacks would require
additional geometric input not supplied in this chapter.
\end{proof}

\begin{remark}[The single-phenomenon principle]
\label{rem:single-phenomenon}
Proposition~\ref{prop:five-theorems-geometry} is not five separate
theorems unified by analogy.  It is a single formal-local principle:
the derived self-intersection admits the standard bar/Koszul model
$\barB(B)$ and its Hochschild shadow is the shifted-cotangent model
(formalized as Theorem~\ref{thm:steinberg-presentation}):
the groupoid structure (A), the expected \'etaleness picture (B), the
$(-1)$-shifted symplectic form on a complementary intersection (C),
the deformation over moduli of curves (D), and the HKR isomorphism
(H).  The Lagrangian self-intersection is the organizing geometry,
and the five theorems are its five faces.
\end{remark}

\begin{remark}[\'Etale groupoid and Koszul Morita theorem]
\label{rem:etale-cross-ref}
Heuristically, for Koszul algebras the self-intersection groupoid
$\Steinb \rightrightarrows \cL$ should be \'etale
\textup{(}Proposition~\textup{\ref{prop:etale-groupoid-koszul}}\textup{)}:
both source and target maps are expected to be quasi-isomorphisms of
derived stacks, because the higher Tor terms
$\operatorname{Tor}^{\cO(\cM)}_i(B,B) = 0$ vanish for $i > 0$.
Koszul inversion is then expected to be the corresponding groupoid
Morita theorem \textup{(}Corollary~\textup{\ref{cor:koszul-morita}}\textup{)}.
Off the Koszul locus, the groupoid acquires genuine derived
fibers---the higher $A_\infty$ operations---and the Morita
equivalence is replaced by completion.  The transition from the
unconditional categorical equivalence of
Theorem~\ref{thm:steinberg-presentation}(ii) to the Koszul
specialisation of~(iii) is the passage from the formal derived
groupoid model to its conjectural \'etale shadow.
\end{remark}

\begin{proposition}[Koszul duality as complementary Lagrangian symmetry;
\ClaimStatusProvedHere]
\label{prop:koszul-serre}
\index{Koszul duality!as complementary Lagrangian symmetry|textbf}
\index{complementary Lagrangians!Koszul involution}
Koszul duality $\cA \leftrightarrow \cA^!$ is the exchange of
complementary Lagrangians in the derived symplectic category.
Precisely:
\begin{enumerate}[label=\textup{(\roman*)}]
\item The Lagrangians $\cL_\cA$ and $\cL_{\cA^!}$ are complementary
  in $\cM$: $T_{\cL_\cA}\cM \simeq T\cL_\cA \oplus T\cL_{\cA^!}$.
\item Their intersection $\cL_\cA \times_\cM^h \cL_{\cA^!}$
  carries the $(-1)$-shifted symplectic form that is the bar-cobar
  pairing
  $\langle\barB(\cA),\, \Omegach(\barB(\cA))\rangle$.
\item The involution $(\text{-})^{!!} \simeq \id$
  \textup{(}double Koszul duality is the identity on the Koszul
  locus\textup{)} follows from the \emph{self-duality} of the
  bar-cobar pairing, which is symmetric because the ambient
  $(-2)$-shifted symplectic form on~$\cM$ is symmetric.
  Geometrically, this is the statement that complementarity is
  symmetric: if\/ $\cL_\cA$ and\/ $\cL_{\cA^!}$ are complementary
  Lagrangians, then\/ $\cL_{\cA^!}$ and\/ $\cL_{(\cA^!)^!}$ are
  also complementary, and the resulting Lagrangian\/
  $\cL_{(\cA^!)^!}$ coincides with\/ $\cL_\cA$.
\end{enumerate}
\end{proposition}

\begin{proof}
(i) is Vol~I, Theorem~C: the deformation and obstruction complexes
of $\cA$ and $\cA^!$ decompose $T\cM|_\cL$ into complementary
Lagrangian summands.  (ii)~follows from PTVV, Theorem~2.9: the
intersection of complementary Lagrangians inherits a $(-1)$-shifted
symplectic form, and this form is the non-degenerate pairing between
bar and cobar.

(iii)~The $(-2)$-shifted symplectic form $\omega_{-2}$ on $\cM$ is
symmetric (even shift), so the bar-cobar pairing
$\langle \barB(\cA), \Omegach(\barB(\cA)) \rangle$ --- which is the
restriction of $\omega_{-2}$ to the Lagrangian intersection --- is
a \emph{self-dual} pairing.  Algebraically: on the Koszul locus, the
bar-cobar adjunction is a Quillen equivalence
(Vol~I, Theorem~B): $\Omegach(\barB(\cA)) \xrightarrow{\sim} \cA$.
Applying the bar construction and then Koszul duality twice gives
$(\cA^!)^! = \Omegach(\barB(\cA^!))$.  Since $\cA^!$ is Koszul when
$\cA$ is (by the symmetry of the complementary Lagrangian
decomposition), the bar-cobar unit for $\cA^!$ is also an
equivalence, so $(\cA^!)^! \simeq \cA$.  Geometrically: the
complementary decomposition $T\cM|_\cL \simeq T\cL_\cA \oplus
T\cL_{\cA^!}$ is symmetric in its two summands (because
$\omega_{-2}$ is symmetric, so the roles of the two Lagrangian
complements are interchangeable), and
exchanging the roles of $\cA$ and $\cA^!$ gives
$T\cM|_{\cL_{\cA^!}} \simeq T\cL_{\cA^!} \oplus T\cL_\cA$, with
the second complement recovering $\cL_\cA$.

\smallskip
\noindent\emph{Clarification.}
The involution $(\text{-})^{!!} \simeq \mathrm{id}$ does \emph{not}
follow from $S^2 = \mathrm{id}$ for a Serre functor.  For a
Calabi--Yau variety of dimension~$n$, the Serre functor is
$S = [n] \otimes \omega$ and $S^2 = [2n] \otimes \omega^2$, which is
generally nontrivial.  The mechanism here is different: the
$(-2)$-shifted symplectic form is symmetric (even degree), so the
induced pairing on the bar-cobar complex is self-dual, and this
self-duality forces $(-)^{!!} \simeq \mathrm{id}$ on the Koszul
locus.
\end{proof}

\begin{remark}[The formal/non-formal dichotomy as the perturbative/non-perturbative divide]
\label{rem:koszul-perturbative-watershed}
\index{Swiss-cheese formality!perturbative/non-perturbative dichotomy|textbf}
\index{perturbative!Lagrangian intersection characterization}
\index{non-perturbative!excess intersection}
The formal/non-formal Swiss-cheese dichotomy is the mathematical form
of the perturbative/non-perturbative divide.  An algebra with formal
Swiss-cheese structure has a \emph{clean} Lagrangian
self-intersection: the derived intersection
$\cL_\cA \times_\cM^h \cL_\cA$ is concentrated in degree zero, the
QFT is perturbatively exact (one-loop, no higher corrections), genus
corrections are scalar multiples of $\omega_g$, and the bar-complex
coproduct is genus-independent (though the $R$-matrix may still
acquire genus dependence through the propagator channel when the
Coisson bracket $c_0 = \{a\,{}_0\,b\}$ is nonzero).
An algebra with non-formal Swiss-cheese structure has an \emph{excess}
intersection: the derived self-intersection carries higher cohomology,
higher $\Ainf$ operations proliferate like instanton corrections,
curvature and braiding entangle
(Theorem~\ref{thm:pole-structure-dichotomy}), and the genus tower
carries genuinely new data at each level.  The OPE pole structure
determines the intersection type: simple poles give transverse meeting
(perturbative); higher-order poles create excess dimension in the
derived fiber (non-perturbative).
(Note: all standard-landscape algebras---including Virasoro and
$\mathcal{W}_N$---are chirally Koszul by PBW universality;
the formal/non-formal distinction classifies complexity
\emph{within} the Koszul locus.)
\end{remark}

%%%%%%%%%%%%%%%%%%%%%%%%%%%%%%%%%%%%%%%%%%%%%%%%%%%%%%%%%%%%%%%%%%%%%%%%%%%%%
% SECTION 1.2 — RELATION TO RAVIOLO VERTEX ALGEBRAS
%%%%%%%%%%%%%%%%%%%%%%%%%%%%%%%%%%%%%%%%%%%%%%%%%%%%%%%%%%%%%%%%%%%%%%%%%%%%%
